\subsection{Composing with the Inverse} 
\begin{fact}[Composition with the Inverse]
For any one-to-one function $f$:
\begin{itemize}
\item For each $x\in\dom f$, $(f\inv\circ f)(x)=x$
\item For each $y\in\rng f$, $(f\circ f\inv)(y)=y$
\end{itemize}
In other worse, composing with the inverse ``undoes'' the function.
\end{fact}
\begin{Example}Let $f(x)=\frac{2}{5}x-4$ and $g(x)=\frac{5}{2}x+10$. Find $f\circ g$ and $g\circ f$ to confirm that $f=g\inv$.
\begin{lpic}{}{7}
\end{lpic}
\end{Example}
\begin{Note}When checking inverses, composing in one order is enough, assuming the domains and ranges align.\end{Note}